\documentclass{report}
\usepackage{amsmath,amssymb}
\setlength{\parindent}{0mm}
\setlength{\parskip}{1em}
\begin{document}
\begin{center}
\rule{6in}{1pt} \
{\large Oliver Lass \\
{\bf Robust optimization using a second order approximation technique in parametrized shape optimization }}

Technische Universit\"at Darmstadt \\ Fachbereich Mathematik \\ Dolivostr 15 \\ 64293 Darmstadt \\ Germany
\\
{\tt lass@mathematik.tu-darmstadt.de}\\
Stefan Ulbrich\end{center}

A nonlinear constrained optimization problem with uncertain parameters is
investigated. It is adressed by a robust worst-case formulation. The
resulting optimization problem is of bi-level structure. This type of
problems are difficult to treat numerically. We propose and investigate
an approximate robust formulation that employs a quadratic approximation.
For an efficient realization in application problems we will mix the
introduced framework with a linear approximation when appropriate.

The proposed strategy is then applied to the optimal placement of a
permanent magnet in the rotor of a synchronous machine. The goal is to
optimize the volume and position while maintaining a desired performance
level. These quantities are computed from the magnetic vector potentials
obtained by the magnetostatic approximation of Maxwell's equation with
transient movement of the rotor. Permanent magnet synchronous machines
can be described sufficiently accurate by this model. Hence we arrive at
an optimization problem governed by a set of elliptic partial
differential equations, where one PDE has to be solved for every rotor
position.

Due to manufacturing, there are uncertainties in material and production
precision. Here the introduced robust optimization framework comes into
play and accounts for uncertain model and optimization parameters. To
speed up the computation in the optimization process reduced order models
are developed. Numerical results are presented to validate the presented
approach.


\end{document}
